% q0001.tex — TEMPLATE / EXAMPLE QUESTION
% Copy this file to add a new question. Rename to q000X.tex.
% Do NOT compile this file directly — use a compiled/*.tex master file.

\question{
  id        = {0001},
  title     = {Elasticity of Demand},
  source    = {OBECON},        % OBECON | IEO | CEO | REO | NET | NOIC
  year      = {2023},
  round     = {Phase 1},       % e.g. Phase 1, Phase 2, Final, Semifinal
  language  = {Portuguese},    % Portuguese | English | Russian | French
  topic     = {Micro},
  subtopic  = {Demand Theory},
  type      = {objective},     % objective | dissertative
  statement = {%
    Se o preço de um bem aumenta de R\$10 para R\$12 e a quantidade demandada
    cai de 100 para 80 unidades, qual é a elasticidade-preço da demanda
    (em valor absoluto)?
  },
  choices   = {%
    \item 0{,}5
    \item 1{,}0
    \item 1{,}5
    \item 2{,}0
    \item 2{,}5
  },
  answer    = {B},
  solution  = {%
    Usando a fórmula da elasticidade ponto:
    \[
      \varepsilon = \left|\frac{\Delta Q / Q}{\Delta P / P}\right|
                  = \left|\frac{-20/100}{+2/10}\right|
                  = \left|\frac{-0{,}2}{0{,}2}\right| = 1{,}0
    \]
    Portanto, a demanda é unitariamente elástica. Alternativa \textbf{(B)}.
  },
}
